\chapter{Theoretical Background}

\section{Probability Theory}

Human driving behaviour is inherently stochastic. In a traffic scenario, different drivers may interpret the same situation differently depending on their perception, 
experience, and intentions. As a result, multiple maneuvers may be possible under the same traffic conditions. For example, a driver approaching a slower vehicle may 
choose to brake, maintain speed, or initiate a lane change. Since such decisions cannot be determined deterministically, behaviour prediction must account for uncertainty.
Probability theory provides a mathematical framework for representing and reasoning about uncertainty. Instead of predicting a single deterministic outcome, probabilistic 
models represent a distribution over possible outcomes and their likelihoods. This allows the model to capture the variability and uncertainty inherent in human decision-making.

\par In probability theory, each uncertain quantity is represented as a \ac{RV} \cite{pml1Book}. 
In the context of driver behaviour modelling, variables such as vehicle velocity, acceleration, or lane position can be treated as observable \acp{RV}. 
In addition, latent quantities such as driver style or driving maneuver intention can also be represented as \acp{RV} that must be inferred from observations.

\par Complex systems require modelling interactions between many \acp{RV}. When several \acp{RV} are considered together, their behaviour can be described using a joint probability distribution \cite{pml1Book}. 
 For two \acp{RV}, $O$ and $Z$, representing the observable \ac{RV} and latent \ac{RV}, the joint probability is written as $P(O,Z)$ and represents the probability that both 
 variables take specific values at the same time. 
 This idea can be extended to any number of variables. In driver behaviour modelling, joint probability is important 
 because observable variables and latent variables influence each other and must often be described together within a unified probabilistic framework.

 \par It is also important to describe dependencies between \acp{RV}. This is done using conditional probability \cite{pml1Book}. The conditional probability of $O$ 
 given $Z$ is defined as
\[
P(O \mid Z) = \frac{P(O,Z)}{P(Z)},
\]
provided that $P(Z) > 0$. Conditional probability describes how the probability of one \ac{RV} changes when information about another \ac{RV} is known. 
This concept is essential because the probability of a future driving maneuver is estimated conditioned on the observed traffic context during inference. 
During model learning, the probability of observations conditioned on latent states is also required.

\par Another important concept is marginal probability, which gives the probability of a subset of \acp{RV} irrespective of the values of the others \cite{pml1Book}. For example, 
if the joint distribution $P(O,Z)$ is known, the marginal probability of $Z$ can be obtained by summing over all possible values of $O$:
\[
P(Z) = \sum_O P(O,Z).
\]

\par In addition to describing dependencies between \acp{RV}, it is also useful to identify situations where variables do not influence each other. 
Two \acp{RV} are said to be independent if the value of one variable does not affect the probability of the other. Formally, two variables $A$ and $B$ are independent 
if their joint probability can be written as
\[
P(A,B) = P(A)P(B).
\]

\par In many practical problems, variables are not completely independent but may become independent when conditioned on another variable. This property is 
known as conditional independence. Two variables $A$ and $B$ are conditionally independent given a third variable $C$ if the joint conditional distribution factorizes as
\[
P(A,B \mid C) = P(A \mid C)P(B \mid C).
\]

Conditional independence is a key concept in probabilistic graphical models because it allows complex joint probability distributions to be represented using simpler 
conditional relationships between variables \cite{koller2009probabilistic}.

\subsection{Bayes' Theorem}

A fundamental principle in probabilistic reasoning is Bayes' theorem, which allows beliefs about uncertain events to be updated when new observations become available \cite{pml1Book}. 
It relates conditional and marginal probabilities and is expressed as

\begin{equation}
P(Z \mid O) = \frac{P(O \mid Z) P(Z)}{P(O)}
\label{eq:bayes}
\end{equation}

\par Here, $P(Z)$ is called the prior probability. The prior represents the initial belief about a variable before observing any new data. In the context of driver 
behaviour modelling, the prior can represent the initial probability of different latent driver states before considering the current traffic situation.
The term $P(O \mid Z)$ is called the likelihood. The likelihood describes how probable the observed data $O$ is, given a particular latent state $Z$. In driver behaviour 
prediction, it measures how compatible the observed vehicle motion and traffic features are with a given latent driver state.
The term $P(O)$ is called the evidence (or marginal likelihood). It represents the total probability of observing the data $O$ under all possible values of $Z$. 
The evidence acts as a normalization factor that ensures the resulting probabilities sum to one.
Finally, $P(Z \mid O)$ is called the posterior probability. The posterior represents the updated belief about $Z$ after observing $O$. In behaviour modelling, 
the posterior corresponds to the updated probability distribution over the latent driver state given the observed traffic context.


\par Bayes' theorem therefore provides a systematic way to update probabilistic beliefs when new information becomes available. 
This principle forms the basis for probabilistic inference in many models used for behaviour prediction, including \ac{DBN}.

\section{Probabilistic Graphical Models}

Probabilistic graphical models provide an efficient way to represent complex joint probability distributions over many \acp{RV} \cite{pml1Book} \cite{koller2009probabilistic}. 
As the number of variables increases, modelling the full joint distribution becomes difficult because the dependencies between variables also increase. 
Probabilistic graphical models address this problem by representing \acp{RV} and their probabilistic relationships using a graph. 
In this representation, nodes correspond to \acp{RV}, while edges describe dependencies between them. 
The graph structure also encodes conditional independence assumptions. These assumptions allow the joint probability distribution to 
be factorised into simpler \acp{CPD}.

\par There are two main types of probabilistic graphical models:
\begin{enumerate}
    \item \textbf{Directed graphical models}, where edges have a direction and represent conditional dependencies between variables.
    \item \textbf{Undirected graphical models}, where edges do not have a direction and represent symmetric relationships between variables.
\end{enumerate}

\par One important class of probabilistic graphical models is the \ac{BN}. Its temporal extension is the \ac{DBN}, which is particularly suitable for modelling 
sequential processes over time.


\subsection{Bayesian Networks}

A \ac{BN} is a directed probabilistic graphical model that represents a set of \acp{RV} and their conditional dependencies using a \ac{DAG} \cite{koller2009probabilistic}. 
Each variable in the network is associated with a \ac{CPD}, which describes how the variable depends on its parent variables. 
These \acp{CPD} are often represented using \acp{CPT}. If a variable has no parents in the graph, its distribution is given by a prior probability distribution. 

\par In general, the joint probability distribution of all variables in a \ac{BN} can be factorised according to the graph structure. For a set of variables 
$X_1, X_2, \dots, X_n$, the joint distribution can be written as
\begin{equation}
P(X_1, X_2, \dots, X_n) = \prod_{i=1}^{n} P(X_i \mid \mathrm{Pa}(X_i))
\label{eq:bn_factorization}
\end{equation}

Here, $\mathrm{Pa}(X_i)$ denotes the set of parent variables of $X_i$ in the graph. This factorisation reflects the conditional independence 
assumption embedded in the network structure.

\begin{figure}[htbp]
    \centering
    \begin{tikzpicture}[
    node distance=2.5cm,
    >=stealth,
    every node/.style={draw, circle, minimum size=1cm, align=center}
]

\node (S) {$S$};
\node[right of=S] (A) {$A$};
\node[below right=1.5cm and 1.25cm of S] (O) {$O$};

\draw[->] (S) -- (A);
\draw[->] (S) -- (O);
\draw[->] (A) -- (O);

\end{tikzpicture}
    \caption{Example of a \ac{BN}}
    \label{fig:bn}
\end{figure}
For the example network shown in \autoref{fig:bn}, the joint distribution can be written as
\[
P(S,A,O) = P(S)P(A \mid S)P(O \mid S,A).
\]

\par Standard \acp{BN} describe relationships between variables within a single time step. However, in many real-world problems, such as 
driver behaviour modelling, the system evolves over time. In such cases, a temporal extension of \acp{BN}, called a \ac{DBN}, 
is used.

\subsection{Dynamic Bayesian Networks}

In a \ac{DBN}, the system is represented as a sequence of time-indexed \acp{RV} that describe how the system evolves over time. 
A common representation of a \ac{DBN} uses a two-slice temporal model that captures dependencies both within a single time step 
and across consecutive time steps. This representation typically relies on the first-order Markov assumption, which states that the current state depends 
only on the previous state and not on earlier states.

\begin{figure}[htbp]
    \centering
    \begin{tikzpicture}[
    node distance=2.5cm,
    >=stealth,
    latent/.style={draw, circle, minimum size=1cm, align=center},
    obs/.style={draw, rectangle, minimum width=1.2cm, minimum height=0.8cm, align=center}
]

\node[latent] (Zt) {$Z_t$};
\node[latent, right=3cm of Zt] (Ztp1) {$Z_{t+1}$};

\node[obs, below=1.8cm of Zt] (Ot) {$O_t$};
\node[obs, below=1.8cm of Ztp1] (Otp1) {$O_{t+1}$};

\draw[->] (Zt) -- (Ztp1);
\draw[->] (Zt) -- (Ot);
\draw[->] (Ztp1) -- (Otp1);

\end{tikzpicture}
    \caption{Two-slice representation of a \ac{DBN}.}
    \label{fig:dbn}
\end{figure}

\par In the example shown in \autoref{fig:dbn}, let $Z_t$ denote the latent state at time step $t$ and $O_t$ denote the corresponding observation. Under the Markov 
assumption, the joint probability distribution over a sequence of states and observations can be written as

\begin{equation}
P(Z_{1:T}, O_{1:T}) =
P(Z_1)
\prod_{t=2}^{T} P(Z_t \mid Z_{t-1})
\prod_{t=1}^{T} P(O_t \mid Z_t).
\label{eq:dbn_factorization}
\end{equation}

\par In this formulation, $P(Z_t \mid Z_{t-1})$ represents the state transition model, which captures how the latent state evolves over time, 
while $P(O_t \mid Z_t)$ represents the emission model that relates the hidden state to the observed data.

\par In the context of driver behaviour modelling, \acp{DBN} provide a probabilistic framework for representing latent driver states and their temporal evolution 
based on observed vehicle motion and traffic context.




\section{Inference in Chain-Structured DBNs}

Inference in probabilistic graphical models refers to computing the probability distribution of unknown variables given observed evidence. 
In temporal models such as \acp{DBN}, this involves estimating the distribution of latent states from a sequence of observations \cite{pml1Book} \cite{koller2009probabilistic}.

\par For the \ac{DBN} shown in \autoref{fig:dbn}, the latent variable $Z_t$ may represent hidden quantities such as driver style or maneuver state. 
These variables are not directly observable and must therefore be inferred from measurable quantities such as vehicle motion and the surrounding traffic context. 
Inference is required because the model contains hidden states, whereas only the observation sequence $O_{1:T}$ is available.

\par In temporal probabilistic models, three common inference tasks arise \cite{Murphy2012MachineL}:
\begin{enumerate}
    \item \textbf{Filtering:} Estimating the current latent state using observations up to the current time step, $P(Z_t \mid O_{1:t})$.
    \item \textbf{Smoothing:} Estimating the latent state using the full observation sequence, $P(Z_t \mid O_{1:T})$.
    \item \textbf{Prediction:} Estimating future latent states from the observations available up to the current time, $P(Z_{t+k} \mid O_{1:t})$.
\end{enumerate}


\par For the chain-structured \ac{DBN} shown in \autoref{fig:dbn}, the latent state at time $t$ depends only on the previous state $Z_{t-1}$, and the observation $O_t$ depends only on the current state $Z_t$. 
This structure corresponds to a first-order Markov model with conditionally independent observations, which is equivalent to the structure of a \ac{HMM}. 
Therefore, exact inference can be performed using the forward-backward algorithm \cite{rabiner1989tutorial}. 
This algorithm is widely used in \acp{HMM}, which are a special case of \acp{DBN}.



\subsection{Forward–Backward Algorithm}

The Forward--Backward algorithm computes posterior probabilities of latent states by combining information from past and future observations \cite{rabiner1989tutorial}. 
More specifically, it estimates the posterior distribution of the hidden state at each time step given the complete observation sequence.

\par Let the observation sequence be \(O_{1:T} = \{O_1, O_2, \dots, O_T\}\) and the corresponding latent states be \(Z_{1:T}\). 
The main objective is to compute the posterior distribution, $P(Z_t \mid O_{1:T})$, which is called the smoothing distribution.

\par The Forward--Backward algorithm performs this computation using two recursive quantities: the forward messages and the backward messages.



\subsubsection{Forward Pass}
The forward variable \(\alpha_t(z_t)\) represents the joint probability of observing the partial sequence up to time \(t\) and being in state \(z_t\):
\begin{equation}
\alpha_t(z_t) = P(O_{1:t}, Z_t = z_t)
\label{eq:forward_def}
\end{equation}

It can be computed recursively as
\begin{equation}
\alpha_t(z_t) =
P(O_t \mid Z_t = z_t)
\sum_{z_{t-1}}
P(Z_t = z_t \mid Z_{t-1} = z_{t-1})
\alpha_{t-1}(z_{t-1})
\label{eq:forward_recursion}
\end{equation}

with initialization
\[
\alpha_1(z_1) = P(Z_1 = z_1)\, P(O_1 \mid Z_1 = z_1)
\]

Here, $P(Z_t = z_t \mid Z_{t-1} = z_{t-1})$ is the transition probability, $P(O_t \mid Z_t = z_t)$ is the observation likelihood, and $P(Z_1 = z_1)$ is the initial state distribution.




\subsubsection{Backward Pass}

The backward variable \(\beta_t(z_t)\) represents the probability of observing the remaining sequence from time \(t+1\) to \(T\) given the current state:
\begin{equation}
\beta_t(z_t) = P(O_{t+1:T} \mid Z_t = z_t)
\label{eq:backward_def}
\end{equation}

It is computed recursively in the reverse direction as
\begin{equation}
\beta_t(z_t) =
\sum_{z_{t+1}}
P(Z_{t+1} = z_{t+1} \mid Z_t = z_t)
P(O_{t+1} \mid Z_{t+1} = z_{t+1})
\beta_{t+1}(z_{t+1})
\label{eq:backward_recursion}
\end{equation}

with initialization
\[
\beta_T(z_T) = 1
\]



\subsubsection{Posterior State Estimation}

Once the forward and backward messages have been computed, the posterior probability of the hidden state at time step \(t\) can be calculated.
This quantity is commonly denoted as $\gamma_t(z_t)$ and represents the probability that the system is in state $z_t$ at time $t$ given the full observation sequence:
\[
\gamma_t(z_t) = P(Z_t = z_t \mid O_{1:T})
\]

Using the forward and backward variables, this can be written as
\begin{equation}
\gamma_t(z_t) = 
\frac{\alpha_t(z_t)\beta_t(z_t)}{P(O_{1:T})}
\label{eq:gamma_posterior}
\end{equation}

where \(P(O_{1:T})\) is the likelihood of the full observation sequence. Using the forward and backward variables, it can be written as
\[
P(O_{1:T}) = \sum_{z_t} \alpha_t(z_t)\beta_t(z_t).
\]
for any time step $t$.

\par The Forward--Backward algorithm therefore combines information from earlier observations through the forward pass and from later observations through the backward pass.


\subsubsection{Posterior Transition Probability}

In addition to the posterior probability of individual states, it is often useful to compute the joint posterior probability of two consecutive states. 
This quantity gives the posterior probability of a transition from state $z_t$ to state $z_{t+1}$ given the full observation sequence:

\[
\xi_t(z_t, z_{t+1}) =
P(Z_t = z_t, Z_{t+1} = z_{t+1} \mid O_{1:T})
\]

Using the forward and backward variables, this can be written as
\begin{equation}
\xi_t(z_t, z_{t+1}) =
\frac{\alpha_t(z_t)\,
P(Z_{t+1}=z_{t+1}\mid Z_t=z_t)\,
P(O_{t+1}\mid Z_{t+1}=z_{t+1})\,
\beta_{t+1}(z_{t+1})}
{P(O_{1:T})}
\label{eq:xi_posterior}
\end{equation}

This posterior transition probability is useful in parameter learning, especially in \ac{EM} based training methods \cite{rabiner1989tutorial}.



\section{Parameter Learning using Maximum Likelihood Estimation}

A probabilistic model contains several unknown parameters that must be learned from data. For a chain-structured temporal model such as the one shown in \autoref{fig:dbn}, 
these parameters can be collectively written as
\[
\theta = \{\pi, A, B\},
\]
where $\pi$ denotes the initial state distribution, $A$ is the state transition matrix, and $B$ represents the parameters of the emission model.

\par If both the observation sequence $O_{1:T}$ and the latent state sequence $Z_{1:T}$ were known, parameter learning could be performed by maximizing the complete-data likelihood
\begin{equation}
\theta^* = \arg\max_{\theta} P(O_{1:T}, Z_{1:T} \mid \theta)
\label{eq:complete_data_mle}
\end{equation}

However, in latent-variable models the state sequence $Z_{1:T}$ is not observed. Therefore, learning must be based only on the observed data, which gives the \ac{MLE} objective
\begin{equation}
\theta^{*} = \arg\max_{\theta} P(O_{1:T} \mid \theta)
\label{eq:observed_data_mle}
\end{equation}

\par Since the latent state sequence is unknown, the likelihood of the observations is obtained by marginalizing over all possible latent state sequences:
\begin{equation}
P(O_{1:T} \mid \theta) = \sum_{Z_{1:T}} P(O_{1:T}, Z_{1:T} \mid \theta)
\label{eq:observed_likelihood_marginalized}
\end{equation}

\par As shown in \eqref{eq:observed_likelihood_marginalized}, the observed-data likelihood is obtained by summing the joint probability over all possible latent state sequences. 
Because this summation involves many possible latent state sequences, direct maximization of the observed-data likelihood in \eqref{eq:observed_data_mle} is difficult. 
To address this problem, the \ac{EM} algorithm is commonly used \cite{Murphy2012MachineL}. 
For chain-structured models such as \acp{HMM}, this learning procedure is known as the Baum--Welch algorithm \cite{rabiner1989tutorial}.




\subsection{Expectation--Maximization Algorithm}

The \ac{EM} algorithm is an iterative method for parameter estimation in probabilistic models with latent variables. 
It addresses the difficulty of maximizing the observed-data likelihood by alternating between estimating the latent states and updating the model parameters.



\subsubsection{Expectation Step}

In the expectation step (E-step), the current parameter estimate is used to compute the posterior distribution of the latent states given the observations. 
For the chain-structured model considered here, this is done using the Forward--Backward algorithm described by \eqref{eq:forward_recursion} and \eqref{eq:backward_recursion}. 
The E-step provides the posterior state probabilities
\[
\gamma_t(i) = P(Z_t = i \mid O_{1:T})
\]
and the posterior transition probabilities
\[
\xi_t(i,j) = P(Z_t = i, Z_{t+1} = j \mid O_{1:T}).
\]


\par As shown in \eqref{eq:gamma_posterior} and \eqref{eq:xi_posterior}, these posterior quantities represent the probabilities of state occupancy and state transitions 
conditioned on the full observation sequence.

\subsubsection{Maximization Step}

In the maximization step (M-step), the model parameters are updated using the posterior quantities computed in the E-step. 

\par For discrete-state chain-structured models, the initial state distribution is updated as
\begin{equation}
\pi_i = \gamma_1(i)
\label{eq:em_initial_update}
\end{equation}

\par The transition probabilities can be re-estimated as
\begin{equation}
A_{ij}
=
\frac{\sum_{t=1}^{T-1} \xi_t(i,j)}
{\sum_{t=1}^{T-1} \gamma_t(i)}
\label{eq:em_transition_update}
\end{equation}


\par The parameters of the emission model are updated using the observations weighted by the posterior state probabilities. 
The exact form of this update depends on the chosen emission model.


\par The E-step and M-step are repeated until convergence. 
Under standard \ac{EM}, each iteration does not decrease the observed-data likelihood, and the algorithm converges to a local optimum or stationary point \cite{Murphy2012MachineL}.